\documentclass[17pt,a4paper,landscape]{extarticle}
\usepackage[margin=2.35cm,top=0.12cm,left=1.05cm]{geometry}
\usepackage{amsmath}
\usepackage{amssymb}
\usepackage{tasks}
\usepackage{tikz}
\usepackage{xcolor}
\usepackage[sfdefault,lf]{FiraSans}
\renewcommand{\familydefault}{\sfdefault}
\setlength{\parindent}{0pt}
\pagestyle{empty}
\settasks{label=\textbf{\arabic*.}, label-width=2.2em, item-indent=2.95em, column-sep=2.1cm, after-item-skip=3.8em}
\renewcommand{\arraystretch}{1.15}
\everymath{\displaystyle}

\begin{document}
\sffamily
\boldmath

\boldmath

\boldmath

\boldmath

\boldmath

\boldmath

\boldmath

\boldmath

\boldmath

\boldmath

\boldmath

\boldmath

\vspace*{0.4em}
{\LARGE \textbf{Excellence}}
\begin{tasks}(3)
\task Explain whether $\frac{14}{21}$ and $\frac{2}{3}$ are equivalent.
\task Explain whether $\frac{18}{24}$ and $\frac{3}{5}$ are equivalent.
\task One fraction is $\frac{3}{4}$. Write three different equivalent fractions.
\task Fill in both blanks: $\frac{\square}{12} = \frac{3}{4} = \frac{21}{\square}$
\task A student says $\frac{4}{6}$ and $\frac{8}{10}$ are equivalent because both numbers doubled. Are they correct?
\task Shade a diagram to show a fraction equivalent to $\frac{2}{5}$.
\task Complete: if $\frac{a}{b} = \frac{12}{20}$, then a simpler equivalent fraction is \rule{2cm}{0.15mm}.
\task Which does not belong: $\frac{1}{2}$, $\frac{2}{4}$, $\frac{3}{6}$, $\frac{4}{10}$?
\task Write a real-life example that could be represented by two equivalent fractions.
\task Write your own pair of equivalent fractions and explain how you know.
\task Explain whether $\frac{15}{20}$ and $\frac{9}{12}$ are equivalent.
\task Fill in both blanks: $\frac{\square}{18} = \frac{2}{3} = \frac{16}{\square}$
\end{tasks}

\end{document}
